\documentclass[11pt,a4paper]{article}

\usepackage[T1]{fontenc}
\usepackage[utf8]{inputenc}
\usepackage{lmodern}
\usepackage{amsmath,amssymb}
\usepackage{geometry}
\usepackage{booktabs}

\geometry{margin=2.5cm}

\begin{document}

\begin{table}[htbp]
\centering
\small
\begin{tabular}{r r r r r r r}
\toprule
$k$ & $p_k$ & $t_k$ & $\lambda_k$ & $p_k^{(\mathrm{out})}$ & $\Delta t$ & $\Delta p$\\
\midrule
1  & 2  & 14.134725 & 14.142881 & 2.001 & -0.008156 & 0.001 \\
2  & 3  & 21.022040 & 21.009774 & 2.999 & -0.012266 & -0.001 \\
3  & 5  & 25.010858 & 25.028331 & 5.003 &  0.017473 & 0.003 \\
4  & 7  & 30.424876 & 30.411209 & 7.001 & -0.013667 & 0.001 \\
5  & 11 & 32.935062 & 32.946800 & 10.998&  0.011738 & -0.002 \\
6  & 13 & 37.586178 & 37.569991 & 13.002& -0.016187 & 0.002 \\
7  & 17 & 40.918719 & 40.903440 & 16.999& -0.015279 & -0.001 \\
8  & 19 & 43.327073 & 43.345912 & 19.004&  0.018839 & 0.004 \\
9  & 23 & 48.005151 & 47.989332 & 23.002& -0.015819 & 0.002 \\
10 & 29 & 49.773832 & 49.781990 & 28.999&  0.008158 & -0.001 \\
\bottomrule
\end{tabular}
\caption{Gemeinsame Realisierung von Zeta-Nullstellen und Primzahlen
durch den arithmetischen Dirac-Operator.
$\Delta t=\lambda_k-t_k$, $\Delta p=p_k^{(\mathrm{out})}-p_k$.}
\label{tab:dirac-first-10}
\end{table}

\end{document}